\documentclass[11pt,a4paper]{article}
% Common preamble for the hint sheets (\input from each file)
\usepackage[utf8]{inputenc}
\usepackage[T1]{fontenc}
\usepackage{lmodern}
\usepackage{amsmath,amssymb,amsthm,mathtools}
\usepackage[a4paper,margin=2.3cm]{geometry}
\usepackage{enumitem}
\usepackage{xcolor}
\usepackage[hidelinks]{hyperref}
\renewcommand{\leq}{\leqslant}
\renewcommand{\le}{\leqslant}
\renewcommand{\geq}{\geqslant}
\renewcommand{\ge}{\geqslant}
\setlength{\parindent}{0pt}
\setlength{\parskip}{0.35em}
% \hintstitle{Sheet title}{Class work / Homework}
\newcommand{\hintstitle}[2]{\begin{center}{\LARGE\bfseries Hints: #1\par}\vspace{0.5em}{\large #2\par}\end{center}
\noindent\textit{The hints to each problem are cumulative. Read one, try again, and open the next only if you are still stuck.}\par\medskip\hrule\medskip}
\newcommand{\problem}[1]{\par\bigskip{\large\bfseries Problem #1}\par\nopagebreak\smallskip}
\newcommand{\subpart}[1]{\par\smallskip\textbf{(#1)}\par\nopagebreak}
\newcommand{\hint}[1]{\par\smallskip\noindent\textcolor{black!70}{\textbf{Hint #1.}}\enspace}
\newcommand{\note}[1]{\par\smallskip\noindent\textit{Note.} #1\par}
\providecommand{\lt}{<}
\providecommand{\gt}{>}
\providecommand{\tightlist}{\setlength{\itemsep}{0pt}\setlength{\parskip}{0pt}}
\emergencystretch=3em

\begin{document}
\hintstitle{Number theory -- 4}{Class work}

\problem{1}
\hint{1} Legendre: $v_p(n!)=\sum_{m\le n}v_p(m)=\sum_{i\ge1}\#\{m\le n:p^i\mid m\}$. For the closed form, write $n=\sum c_jp^j$; then $\lfloor n/p^i\rfloor=\sum_{j\ge i}c_jp^{j-i}$. Sum over $i$.
\hint{2} Kummer: $v_p\binom nk=\frac{s_p(k)+s_p(n-k)-s_p(n)}{p-1}$, and each carry lowers the digit sum by $p-1$.
\hint{3} All $\binom nk$ are odd if and only if adding $k$ and $n-k$ in base $2$ never carries, for every $k$. That means $n$ has no $0$ digit below its top digit: $n=2^m-1$. Next, $v_2\binom{2n}n=s_2(n)$ (doubling carries once per $1$-digit). Finally, $\binom{2n}{n}-\binom{2n}{n+1}=\frac{1}{n+1}\binom{2n}n$ is an integer (Catalan).

\problem{2}
\note{As printed, $v_7(2^{2024}-2)=v_7(2^{2023}-1)=0$, since $\mathrm{ord}_7(2)=3\nmid2023$. So LTE plays no role there. The hints treat the statement as printed.}
\hint{1} LTE, odd $p$: $\frac{a^n-b^n}{a-b}=\sum a^{n-1-i}b^i\equiv na^{n-1}\pmod p$. This proves the case $p\nmid n$. For $n=p$, show that $p\,\|\,\frac{a^p-b^p}{a-b}$ by writing $b=a+tp$. Induct. For $p=2$ and even $n$, factor $a^n-b^n=(a^{n/2}-b^{n/2})(a^{n/2}+b^{n/2})$, or use $a^2-b^2$ first.
\hint{2} $v_2(3^n-1)$ is $1$ for odd $n$ and $v_2(2)+v_2(4)+v_2(n)-1=2+v_2(n)$ for even $n$. Solving $2+v_2(n)\ge n$: answer $n\in\{1,2,4\}$.
\hint{3} $2^{2024}-2=2(2^{2023}-1)$ and $2^{2023}\equiv2\pmod 7$, so the answer is $0$. (For comparison: $v_7(2^{3m}-1)=v_7(8-1)+v_7(m)=1+v_7(m)$.)

\problem{3}
\hint{1} Compare $v_p$ of both sides. Note $v_p(\mathrm{lcm}(1,\dots,m))=\#\{j\ge1:p^j\le m\}$.
\hint{2} $\sum_i\#\{j:p^j\le n/i\}=\sum_j\#\{i:i\le n/p^j\}=\sum_j\lfloor n/p^j\rfloor=v_p(n!)$.

\problem{4}
\hint{1} Every prime factor of $n!+1$ exceeds $n$. Compare with the multinomial coefficient $\frac{(2012n)!}{(n!)^{2012}}$.
\hint{2} $\gcd(n!+1,n!)=1$, so $n!+1$ divides $\frac{(2012n)!}{(n!)^{2012}}\le2012^{2012n}$. But $n!$ grows faster than any $C^n$. Answer: the set is finite.

\problem{5}
\hint{1} Bound $a_n/n$. If $n=ab$ with $1<a<b$, then $n\mid b!$, so $a_n\le b=n/a\le n/2$.
\hint{2} If $n=q^2$ with $q$ prime, then $a_n=2q\le n/2$ for $q\ge5$. The exceptions are $n=4$ ($a_4=4$) and $n=9$ ($a_9=6$).
\hint{3} So apart from two terms, $(a_n/n)^n\le2^{-n}$. Answer: the series converges.
\end{document}
